% 核心结论：作业二：节点守恒与相量闭合
% 图原型：纯几何／对比；按本题相位与电路拓扑绘制。
\documentclass[tikz,border=7pt]{standalone}
\usepackage{fontspec,xeCJK,amsmath,amssymb}
\IfFontExistsTF{Times New Roman}{\setmainfont{Times New Roman}}{\setmainfont{TeX Gyre Termes}}
\IfFontExistsTF{Heiti SC}{\setCJKmainfont[AutoFakeBold=1.5]{Heiti SC}}{\setCJKmainfont[AutoFakeBold=1.5]{FandolHei-Regular}}
\usetikzlibrary{arrows.meta,calc,angles,quotes}
\definecolor{ink}{HTML}{243745}
\definecolor{blue}{HTML}{275DAB}
\definecolor{green}{HTML}{177C63}
\definecolor{red}{HTML}{BD4935}
\definecolor{purple}{HTML}{715393}
\begin{document}
\begin{tikzpicture}[font=\fontsize{9.5}{12}\selectfont,text=ink,
 v/.style={line width=1.1pt,-{Stealth[length=2.4mm,width=1.5mm]}},
 axis/.style={ink!65,line width=.45pt,-{Stealth[length=1.7mm,width=1.1mm]}},
 aux/.style={ink!40,densely dashed,line width=.5pt},
 lab/.style={fill=white,inner sep=2pt},
 card/.style={draw=blue!25,fill=blue!3,rounded corners=3pt,text width=6cm,align=left,inner sep=10pt}]
\path[use as bounding box] (-8,-3.35) rectangle (8,3.35);
\node[font=\fontsize{11}{14}\selectfont\bfseries] at (0.00000,3.04000) {作业二：节点守恒与相量闭合};
\begin{scope}[shift={(-4,0)}]
\draw[v,blue] (-2.00000,1.25000) -- (0.00000,0.00000);
\node[lab,text=blue] at (-1.50000,1.35000) {$i_1$};
\draw[v,green] (-2.00000,-1.25000) -- (0.00000,0.00000);
\node[lab,text=green] at (-1.50000,-1.35000) {$i_2$};
\draw[v,red] (0.00000,0.00000) -- (2.10000,0.00000);
\node[lab,text=red] at (1.30000,0.35000) {$i_3$};
\fill[ink] (0,0) circle (2pt);
\node[] at (0.05000,0.45000) {$N$};
\node[] at (0.00000,-2.35000) {$i_1+i_2-i_3=0$};
\end{scope}
\begin{scope}[shift={(4,0)}]
\draw[v,blue] (-1.80000,-1.25000) -- (1.20000,-1.25000);
\node[lab,text=blue,anchor=center] at (-0.30000,-1.60000) {$4$};
\draw[v,green] (1.20000,-1.25000) -- (1.20000,1.00000);
\node[lab,text=green,anchor=west] at (1.45000,-0.12500) {$j3$};
\draw[v,red] (-1.80000,-1.25000) -- (1.20000,1.00000);
\node[lab,text=red,anchor=east] at (-0.45000,0.02500) {$\widetilde I_3$};
\node[] at (-2.00000,-1.37000) {$O$};
\node[] at (1.40000,-1.45000) {$A$};
\node[] at (1.40000,1.20000) {$B$};
\node[] at (0.00000,-2.35000) {$\widetilde I_3=4+j3\;\mathrm{A}$};
\end{scope}
\node[text width=15cm,align=center] at (0.00000,-3.00000) {左：实际节点及参考方向。右：复平面中的电流相量，不能把两类方向混同。};
\end{tikzpicture}
\end{document}
